\section{Introduction}
\label{sec:introduction}

The bulk of automated vulnerability-detection research over the past
five years has focused on C and C++. Recent benchmark and dataset
efforts --- BigVul~\cite{TODO_bigvul}, CVEfixes~\cite{TODO_cvefixes},
CrossVul~\cite{TODO_crossvul}, DiverseVul~\cite{TODO_diversevul}, and
the recent CASTLE benchmark~\cite{TODO_castle} --- concentrate on
memory-safety bugs, buffer overflows, and use-after-free patterns that
dominate C-language CVE feeds. Python receives substantially less
attention despite being the dominant language for security-critical
infrastructure: machine-learning pipelines, scientific computing, web
back-ends, automation tooling, and cloud-native applications. The
vulnerability classes that matter for Python software --- SQL
injection through ORM and DB-API patterns, deserialization of pickle
and YAML data, code injection via \texttt{eval}/\texttt{exec},
server-side request forgery in webhook endpoints --- are
sink-anchored in ways the C-centric literature does not address.

When researchers do build Python vulnerability datasets, they
typically inherit labels directly from CVE fix commits --- a strategy
known to produce 25--50\% noise from co-changed
files~\cite{TODO_croft}. The labeled ``vulnerable'' file in a CVE fix
is not always the actual sink site: commit boundaries cut across many
files, and CVE descriptions rarely identify the precise line where the
vulnerability lives. Recent audits of BigVul~\cite{TODO_croft}
quantified this at $\approx 25\%$; informal audits of CVE-fix-derived
data we conducted as part of this work found 52\% mislabeled positives
in the raw, unfiltered CVE-feed corpus. Training and evaluation on
this substrate produces models and tool comparisons that partially
measure noise rather than detection capability.

The CASTLE benchmark~\cite{TODO_castle} addressed this evaluation
problem for C/C++ with a hand-curated 250-program corpus and a
comparative study of fourteen static analyzers and ten large language
models. CASTLE accepts CVE/advisory labels at face value but
compensates with small, carefully selected samples. For Python, no
comparable benchmark exists --- researchers fall back to either using
CVE-fix-derived datasets with known noise, or using
Juliet-style~\cite{TODO_juliet} synthetic samples that lack
realistic framework idiom and surrounding application context. Both
substrates limit what can be concluded about tool performance under
real-world conditions. Additionally, CASTLE measures detection rates
but does not evaluate robustness against semantics-preserving
adversarial transformations, leaving open the question of whether
high-scoring tools are doing genuine semantic analysis or pattern
matching that breaks under simple rewriting.

This paper introduces a Python vulnerability-detection benchmark
covering seven sink-shaped MITRE Top-25 CWE classes (CWE-89, CWE-79,
CWE-22, CWE-78, CWE-94, CWE-918, and CWE-502) plus a ``safe''
hard-negative class. We address the label-noise problem with a
three-layer validation methodology that progressively reduces the
audited false-positive rate from 52\% to 26\% --- comparable to
manually audited security datasets at a fraction of the human-hour
cost. We additionally extend CASTLE's evaluation framing in two
directions: four semantics-preserving adversarial mutators isolate
specific detection shortcuts, and per-framework slicing (Flask /
Django / FastAPI) exposes framework-specific gaps in static-analysis
rule coverage. The combined benchmark, label-validation methodology,
and adversarial-robustness probe constitute, to our knowledge, the
first Python evaluation substrate with explicit, automated
label-validity criteria.

\paragraph{Contributions}
\begin{itemize}
  \item \textbf{An audit-validated Python vulnerability benchmark.}
    1{,}300 samples (871 positives plus 429 safe hard negatives) across
    seven sink-shaped MITRE Top-25 CWE classes, with repo-group-stratified
    train / validation / test splits (607 / 127 / 132) and zero
    cross-split repository leakage. Distributed with full per-sample
    metadata, content-hash deduplication, and rejection manifests for
    every dropped sample.

  \item \textbf{A three-layer label-validation methodology.} A pre-ingest
    sink-presence filter, a three-round adversarial audit by independent
    reviewers, and a diff-based filter that requires the sink line to be
    modified between \texttt{code\_before} and \texttt{code\_after}.
    Combined, the three layers reduce audited false-positive rate from
    52\% to 26\%. To our knowledge this is the first CVE-derived Python
    benchmark with explicit, automated label-validity criteria; the
    methodology is reusable across CVE-derived datasets in any language.

  \item \textbf{Four semantics-preserving adversarial mutators}
    (dead-code injection, string splitting, identifier rename, wrapper
    extraction) with deterministic per-sample reproducibility and
    AST-level correctness guarantees. Each mutator targets a known
    detection shortcut. By measuring per-tool F1 drop per mutator, we
    surface which shortcut each tool depends on --- a finer-grained
    diagnosis than CASTLE's aggregate detection rates allow.

  \item \textbf{A comparative evaluation} of $N$ static analyzers
    (Bandit, Semgrep, CodeQL, $\ldots$) and $M$ LLM-based detectors
    (Claude, GPT-4o, $\ldots$) across the benchmark and its
    mutator-derived adversarial variants, reporting per-CWE F1, per-tool
    robustness drop, and per-framework breakdowns.
    % TODO: replace once eval is run with the headline finding —
    % e.g., ``We find that LLMs achieve $\Delta$ F1 advantage over
    % the best SAST baseline on clean samples, but lose $\Delta$
    % F1 points on string-split adversarials, indicating they
    % are not immune to surface-token reliance.''

  \item \textbf{A scope justification} documenting which Top-25 CWEs
    admit sink-anchored detection and which are structural
    (absence-of-code) patterns that require fundamentally different
    methodology. We argue for the explicit separation as a
    contribution: sink-shaped CWEs admit pattern-based evaluation;
    structural CWEs do not, and benchmarks that mix the two confound
    detector capabilities.
\end{itemize}

The remainder of this paper is organized as follows.
Section~\ref{sec:related} surveys vulnerability datasets, CVE-derived
label noise studies, and the static-analysis / LLM detector landscape.
Section~\ref{sec:dataset} describes dataset construction, sources, the
sink-presence filter, the audit methodology, and the diff-based filter
that together define the benchmark's labels.
Section~\ref{sec:methodology} formalizes the evaluation methodology and
the four adversarial mutators.
Section~\ref{sec:evaluation} reports the comparative evaluation
results.
Section~\ref{sec:threats} discusses threats to validity, including the
remaining 26\% audit residual.
Section~\ref{sec:conclusion} summarizes findings and outlines future
work, including the deferred structural-CWE evaluation track and
extension to additional languages.
